%% ma: L12-B16-0012
%% lop: 12
%% chuong: 5
%% bai: 16
%% dang: 12.16.9
%% loai: TN4
%% muc: TH
%% dap_an: "A"
%% nguon: "Ảnh giáo viên gửi 04/11/2026 (Huỳnh Văn Ánh, Chuyên đề V - Phương pháp tọa độ trong không gian), Đề test theo bài, Phần I Câu 3"
%% kien_thuc: "Bài 16 (góc giữa hai mặt phẳng), Bài 14 (vectơ pháp tuyến của mặt phẳng)"
%% trang_thai: chinh-thuc
%% ghi_chu: "Đáp án gốc A khớp (sympy). Đã duyệt 04/11/2026."
\begin{ex}
Trong không gian với hệ trục tọa độ $Oxyz$, cho hai mặt phẳng $(P)\colon 2x-y-z-3=0$ và $(Q)\colon x-z-2=0$. Góc giữa hai mặt phẳng $(P)$ và $(Q)$ bằng
\choice{\True $30^\circ$}{$45^\circ$}{$60^\circ$}{$90^\circ$}
\loigiai{$(P)$ có vectơ pháp tuyến $\vec n_1=(2;-1;-1)$, $(Q)$ có vectơ pháp tuyến $\vec n_2=(1;0;-1)$.

Khi đó $\cos((P),(Q))=\dfrac{|\vec n_1\cdot\vec n_2|}{|\vec n_1|\cdot|\vec n_2|}=\dfrac{|2+0+1|}{\sqrt6\cdot\sqrt2}=\dfrac{\sqrt3}{2}$. Do đó $((P),(Q))=30^\circ$. Chọn A.}
\end{ex}
